% 当前正文配置：依据 TP-0.2；修改需同步受影响的接口与图表。
\newcommand{\PendingValue}{待补充}

% 文档基本信息；PDF 不显示修订编号、稿件状态或封面日期。
\newcommand{\DocumentTitle}{解决方案}
\newcommand{\DocumentAuthor}{\PendingValue}
\newcommand{\DocumentReviewer}{\PendingValue}

% 已知比赛和企业命题信息；学校与具体子方向仍由团队核定。
\newcommand{\CompetitionName}{中国国际大学生创新大赛（2026）}
\newcommand{\CompetitionTrack}{产业赛道}
\newcommand{\CompetitionGroup}{企业命题组}
\newcommand{\CompetitionDirection}{嵌入式智能计算}
\newcommand{\SchoolName}{重庆邮电大学}
\newcommand{\CollegeName}{计算机学院}
% 题名分段组合，封面按三行显示，正文使用完整题名宏。
\newcommand{\ChallengeTitleSoftwarePart}{基于国产开源软件与}
\newcommand{\ChallengeTitleArchitecturePart}{开源RISC-V处理器架构的}
\newcommand{\ChallengeTitleLineOne}{\ChallengeTitleSoftwarePart\ChallengeTitleArchitecturePart}
\newcommand{\ChallengeTitleLineTwo}{嵌入式智能计算系统研究与设计}
\newcommand{\ChallengeTitle}{\ChallengeTitleLineOne\ChallengeTitleLineTwo}
\newcommand{\ChallengeCompany}{中智讯（武汉）科技有限公司}
\newcommand{\ChallengeIdentifier}{239}
\newcommand{\ProjectName}{嵌入式开源智能推理系统}
\newcommand{\ProjectSubtitle}{统一算子开发与异构推理执行}
\newcommand{\ApplicationScenario}{嵌入式模型推理，逐步扩展视觉与智能感知}
\newcommand{\SubmissionStage}{校内选拔赛材料准备}
\newcommand{\SubmissionDeadline}{2026年9月13日12时}

\newcommand{\TechnicalBaselineDate}{2026年9月9日}

\newcommand{\CompanyFeedback}{待企业对接后补充正式意见}

\newcommand{\ProjectKeywords}{嵌入式推理；RISC-V；异构计算；开源硬件；算子开发}
